% ==============================================================================
% Standalone literature review supplement:
%   "Swept Volume Approximation, Interval Kinematics, and Their Prior Use in
%    Endpoint-to-Link Envelope Pipelines"
%
% Purpose: supports the IEEE TRO reviewer response to the novelty concern raised
%   in the review.  This file (a) surveys the classical geometric ingredients
%   individually, (b) documents the closest prior work (Redon et al. 2004/2005,
%   Zhang et al. 2007) that ALREADY implements the two-step endpoint-interval-
%   to-LSS-swept-volume pipeline for CCD of articulated bodies, (c) draws an
%   explicit boundary between what is prior art and what is new in RBF, and
%   (d) provides corrected replacement text for the Related Work section of
%   sbf_tro_2026.tex.
%
% Compile standalone with XeLaTeX:
%   xelatex related_work_envelopes.tex
% ==============================================================================

\documentclass[11pt,a4paper]{article}
\usepackage{fontspec}
\setmainfont{Latin Modern Roman}[Ligatures=TeX]
\usepackage{amsmath,amssymb}
\usepackage{hyperref}
\usepackage{cleveref}
\usepackage{booktabs}
\usepackage[margin=2.5cm]{geometry}
\usepackage[inline]{enumitem}

% Reuse the same shorthand macros as the main paper
\newcommand{\rbf}{\textsc{RBF}}
\newcommand{\etal}{\textit{et al.}}
\newcommand{\Real}{\ensuremath{\mathbb{R}}}

\title{%
  Literature Review Supplement:\\
  Swept Volume Approximation, Interval Forward Kinematics,\\
  and the Endpoint-to-Link Envelope Pipeline\\[0.4em]
  \large Supporting Material for the RapidBoxForest (RBF) TRO Submission}
\author{Yu Tian and Hongliang Ren\\
  Department of Electronic Engineering, CUHK}
\date{May 2026}

\begin{document}
\maketitle
\tableofcontents
\bigskip

% ===========================================================================
\section{Purpose and Scope}
\label{sec:purpose}
% ===========================================================================

The IEEE TRO reviewer raised the concern that the RBF paper claims to be the
``first systematic application'' of Minkowski sum decomposition and convex-hull
generator bounds to compute link interval envelopes for manipulator planning,
but that (a) Minkowski decomposition of capsules is a classical result used in
PQP, GJK, and FCL; (b) convex-hull generator bounding of swept volumes has
ample precedent; and (c) the paper does not cite swept-volume or CCD literature
that already implements the two-step pipeline of bounding link endpoints with
interval arithmetic and then lifting to a link swept volume.

This supplement:
\begin{enumerate}[nosep]
  \item surveys each ingredient individually and its classical status
    (\cref{sec:minkowski,sec:hull_generator,sec:interval_kin});
  \item identifies the closest prior work that has already used the combined
    endpoint-interval-to-swept-capsule pipeline in the CCD context
    (\cref{sec:closest_prior});
  \item draws an explicit boundary between what is prior art and what is novel
    in RBF (\cref{sec:novelty_boundary}); and
  \item provides corrected replacement text for the Related Work section of the
    main paper (\cref{sec:replacement_text}).
\end{enumerate}

% ===========================================================================
\section{Ingredient 1: Minkowski Sum Decomposition of Capsules}
\label{sec:minkowski}
% ===========================================================================

A capsule (line-swept sphere, LSS) is a Minkowski sum $C = S \oplus B_2(r)$,
where $S = [a,b]$ is the core segment and $B_2(r)$ is a Euclidean ball of
radius $r$.  The identity
\begin{equation}
  \mathcal{S}_I(C) = \mathcal{S}_I(S) \oplus B_2(r),
  \label{eq:capsule_mink}
\end{equation}
where $\mathcal{S}_I$ denotes the swept set over any parameter domain $I$,
follows from the fact that the Euclidean ball is isotropic and its summation
commutes with rigid motions.  This is a classical, elementary fact.

\paragraph{Where it appears.}
The earliest systematic exploitation of this decomposition in robotics
collision detection is the Proximity Query Package (PQP) of Larsen
\etal\ \cite{larsen1999fast}, which decomposes every convex primitive into a
``swept sphere volume'' (core shape $\oplus$ radius ball) and converts all
distance and intersection queries to queries on the core shape offset by a
scalar radius.  GJK \cite{gilbert1988fast} and its improved implementation by
van den Bergen \cite{bergen1999gjk} use exactly the same observation: computing
distance from a query point to a capsule reduces to computing distance to the
core segment and subtracting the radius.  FCL \cite{pan2012fcl} generalises
the same interface to a library of swept-sphere primitives.  OBB-tree
construction for capsule-approximated meshes also relies implicitly on
\cref{eq:capsule_mink} when padding OBBs by the link radius
\cite{gottschalk1996obbtree}.

\paragraph{Status.}
\textbf{Classical.}  The identity~\eqref{eq:capsule_mink} is not claimed as
novel by RBF.  It is a foundational tool shared with virtually every
swept-sphere collision library.

% ===========================================================================
\section{Ingredient 2: Convex-Hull Generator Bounding of Swept Volumes}
\label{sec:hull_generator}
% ===========================================================================

For a convex polytope $P = \operatorname{conv}(v_1,\dots,v_m)$ under a family
of rigid transforms indexed by a parameter $\theta \in I$, bounding each
vertex trajectory by a conservative outer set $\widehat{\Gamma}_i(I)$ suffices
to bound the swept polytope:
\begin{equation}
  \mathcal{S}_I(P)
  \subseteq
  \operatorname{conv}\!\Bigl(\bigcup_{i=1}^{m} \widehat{\Gamma}_i(I)\Bigr).
  \label{eq:generator_bound}
\end{equation}
For a capsule link the relevant generators are the two segment endpoints $a$
and $b$, reducing the bounding task to two endpoint trajectory enclosures.

\paragraph{Where it appears.}
The swept-volume literature has used generator bounds for decades.  Early work
by Blackmore and Leu \cite{blackmore1992swept} analyses swept volumes of
convex bodies using Lie groups and establishes that the swept set is
determined by the boundary envelope.  Abdel-Malek \etal\
\cite{abdel2004swept} survey practical swept-volume methods and explicitly note
that for convex bodies, bounding the generator trajectories (vertices of the
convex hull) gives a conservative swept volume approximation.  In the
bounding-volume hierarchy literature, OBB fitting to the swept keyframes of a
convex mesh likewise relies on generator-bounding: the OBB is fitted to the
union of vertex trajectories, not to the full swept-body
manifold \cite{gottschalk1996obbtree}.

In the context of \emph{capsule} links specifically, because the core segment
$[a,b]$ is itself the convex hull of two points (the endpoints), the generator
set reduces to exactly two points.  This reduction is the starting observation
in several continuous collision detection (CCD) works for articulated bodies,
as discussed in \cref{sec:closest_prior}.

\paragraph{Status.}
\textbf{Classical.}  Generator bounding of convex swept sets, and its
reduction to two-endpoint bounding for capsule links, is well-precedented in
both swept-volume theory and CCD practice.  RBF does not claim novelty here.

% ===========================================================================
\section{Ingredient 3: Interval Forward Kinematics for Endpoint Enclosure}
\label{sec:interval_kin}
% ===========================================================================

Interval arithmetic, introduced by Moore \cite{moore1966interval}, provides a
systematic mechanism for computing guaranteed outer bounds on real-valued
functions evaluated over parameter intervals.  The key property is
\emph{inclusion monotonicity}: if inputs are enclosed in intervals, the
interval-arithmetic evaluation of any smooth function encloses the true image.
Applied to kinematics, the idea propagates interval-valued Denavit--Hartenberg
(DH) transforms along the kinematic chain and reads axis-aligned bounding boxes
(AABBs) of endpoint coordinates from the resulting interval-valued prefix
transforms.

\paragraph{Classical applications in robotics.}
Merlet pioneered the use of interval DH-chain propagation for certified
workspace analysis of parallel manipulators: he showed that interval FK can
provide guaranteed bounds on the reachable workspace and detect
singularities \cite{merlet2004solving,merlet2009interval}.  The same
propagation structure---compose interval transforms from root to leaf, read
coordinate bounds at each frame---is used in serial manipulator workspace
analysis, with the known caveat that axis-aligned interval arithmetic over-
approximates the image of composed rotations (wrapping effect).

In computer graphics, Snyder \cite{snyder1992interval} applied interval
arithmetic to parametric geometry: by bounding function coefficients by
intervals, one obtains conservative bounding boxes for parametric curves and
surfaces without mesh sampling.  This established the pattern of using interval
arithmetic as a conservative \emph{spatial filter} rather than as an exact
computation.

\paragraph{Status.}
\textbf{Classical.}  Interval FK for endpoint bounding is an established
kinematic-chain technique.  Its primary design goals in prior work are static
workspace analysis (Merlet) and parametric-curve bounding (Snyder), not
adaptive C-space box growth.

% ===========================================================================
\section{The Closest Prior Work: Redon et al.\ (2004, 2005) and Zhang et al.\ (2007)}
\label{sec:closest_prior}
% ===========================================================================

The most direct prior use of the combined endpoint-interval-to-swept-capsule
pipeline for robot links is the continuous collision detection work of Redon
and collaborators.

\subsection{Redon et al.\ (2004): CCD for Avatars}

In ``Interactive and continuous collision detection for avatars in virtual
environments'' \cite{redon2004continuous_avatar}, Redon, Kim, Lin, and Manocha
model each avatar limb as an LSS (capsule) and compute interval vectors that
bound the coordinates of the LSS \textbf{endpoints} over a given motion
interval.  The two interval-bounded endpoints are then used to derive a
conservative bounding volume for the swept LSS over the motion step.  The
paper explicitly states (paraphrased): \emph{``for each link represented as an
LSS with endpoints $p$ and $q$, we obtain two three-dimensional interval vectors
$[\mathbf{p}]$ and $[\mathbf{q}]$ that bound the positions of the endpoints
over the time interval; this gives a conservative swept-sphere bound for the
link.''}

This is structurally identical to the endpoint-to-link envelope map in RBF:
\begin{enumerate}[nosep]
  \item Compute interval bounds $[\mathbf{p}]$, $[\mathbf{q}]$ on the two
        capsule endpoints over a parameter domain.
  \item Form a conservative link swept-volume bound from
        $\operatorname{bbox}([\mathbf{p}] \cup [\mathbf{q}]) \oplus B_2(r)$.
\end{enumerate}
The difference is that Redon et al.\ use a \emph{time} interval (the motion
step of a given trajectory), whereas RBF uses a \emph{joint-space box} (a
region of C-space for multi-query planning).

\subsection{Redon et al.\ (2005): CCD for General Articulated Models}

The journal version ``Fast continuous collision detection for articulated
models'' \cite{redon2005fast} extends the same framework to arbitrary robot
arms and kinematic trees.  The paper propagates interval transforms along the
kinematic chain to bound each link endpoint, then uses the resulting endpoint
interval vectors to derive conservative bounding capsules for each link's
swept volume.  The algorithm explicitly:
\begin{enumerate}[nosep]
  \item propagates interval DH transforms from root to leaf (interval FK);
  \item at each link frame, reads interval-valued endpoint positions
        $[\mathbf{p}_i]$, $[\mathbf{p}_{i+1}]$;
  \item bounds the swept capsule by
        $\operatorname{bbox}([\mathbf{p}_i] \cup [\mathbf{p}_{i+1}]) \oplus
        B_2(r_i)$, which is then tested against the obstacle.
\end{enumerate}
This constitutes the same two-step pipeline (interval FK endpoint bounding
$\rightarrow$ Minkowski radius lift) that RBF describes as its geometric
contribution.  The paper has accumulated over 240 citations and is a standard
CCD reference.

\subsection{Zhang et al.\ (2007): Taylor-Model CCD for Articulated Bodies}

Zhang, Redon, Lee, and Kim \cite{zhang2007taylor} improve the tightness of the
CCD bounding step by replacing interval arithmetic with \emph{Taylor models},
which track first-order dependency information and substantially reduce wrapping.
The overall pipeline remains the same: bound the motion of each link endpoint
over a time interval, then lift to a swept-capsule bound.  This is analogous to
the hierarchy of endpoint sources (IFK, CritSample, analytical extrema) in
RBF, where IFK corresponds to plain interval arithmetic, and tighter sources
correspond to the Taylor-model or analytical-polynomial improvements.

\subsection{Dynamic Sphere-Swept Line BVs: Corrales et al.\ (2011)}

Corrales, Candelas, and Torres \cite{corrales2011dynamic} apply sphere-swept
line (SSL) bounding volumes dynamically to robot--human interaction safety
monitoring.  Their method tracks the swept capsule of each robot link over a
discrete time step via forward kinematics evaluated at both step endpoints and
uses the resulting SSL envelope as a safety distance margin.  While this work
does not use interval arithmetic for the endpoint bounding, it relies on the
same structural observation: bounding a capsule's motion over a time step
reduces to bounding its two endpoint positions.

\subsection{Adaptive Dynamic Collision Checking: Schwarzer et al.\ (2005)}

Schwarzer, Saha, and Latombe \cite{schwarzer2005adaptive} present an adaptive
collision checking method for articulated robots that, like RBF, organises
collision evidence for reuse across related planning calls.  Their method bounds
the swept volume of each link along a path segment and reuses that bound to
avoid redundant checks.  Endpoint-level bounding (via the robot FK evaluated at
path endpoints) is implicit in their link-motion bound.  Their multi-query reuse
motivation is also similar to RBF, though the prior work addresses path-segment
reuse in a single scene rather than C-space box validation for multi-query
replanning.

% ===========================================================================
\section{Novelty Boundary: Prior Art vs.\ RBF}
\label{sec:novelty_boundary}
% ===========================================================================

\Cref{tab:novelty} summarises the landscape.  The two-step
\emph{endpoint-interval-to-swept-capsule} pipeline is unambiguously present in
Redon et al.\ (2004, 2005) for the CCD purpose.  What differs in RBF is the
\emph{planning purpose and design organisation}:

\begin{table}[htbp]
\centering
\caption{Comparison of prior endpoint-interval-to-link-sweep pipelines vs.\ RBF.}
\label{tab:novelty}
\small
\begin{tabular}{@{}p{3.2cm}p{3.0cm}p{3.0cm}p{4.2cm}@{}}
\toprule
\textbf{Work} & \textbf{Parameter domain} & \textbf{Use case} & \textbf{What is new vs.\ RBF} \\
\midrule
Redon et al.\ 2004/2005
  \cite{redon2004continuous_avatar,redon2005fast}
& Time interval (given path)
& CCD: does the link collide during a single motion step?
& RBF applies over a C-space box; organises as reusable filter \\
\addlinespace
Zhang et al.\ 2007
  \cite{zhang2007taylor}
& Time interval (given path)
& Tighter CCD via Taylor models
& RBF provides three switchable sources (IFK, CritSample, analytical) \\
\addlinespace
Schwarzer et al.\ 2005
  \cite{schwarzer2005adaptive}
& Path segment endpoints
& Adaptive multi-query path reuse
& RBF uses a volumetric box rather than path-level reuse \\
\addlinespace
Corrales et al.\ 2011
  \cite{corrales2011dynamic}
& Discrete time step
& Dynamic human--robot safety
& RBF targets C-space planning, not real-time proximity \\
\midrule
\textbf{RBF (this work)}
& C-space joint-box interval
& Adaptive box-growth and multi-query planning
& Multi-source interface; cached LECT reuse across builds; box-forest planner\\
\bottomrule
\end{tabular}
\end{table}

The \textbf{genuine novelties} in RBF that are not present in any of the above works are:

\begin{enumerate}
  \item \textbf{Application to C-space box validation, not path-time CCD.}
        Prior CCD work bounds link motion over a \emph{given trajectory's time
        step}.  RBF applies the same endpoint-interval bounding to a \emph{C-space
        joint-box interval}: the parameter domain is a joint-limit box, not a
        scalar time step, and the goal is to classify whether any configuration
        in the box is collision-free, not whether a given path collides at some
        time.  This change of domain drives all the downstream design choices
        (depth-synchronous tree growth, LECT, box-forest planner).

  \item \textbf{Multi-source endpoint-envelope interface.}
        Prior CCD work uses a fixed kinematic bounding method (interval
        arithmetic in Redon; Taylor models in Zhang).  RBF explicitly separates
        the endpoint-source choice from the link-representation choice, exposing
        three interchangeable sources (IFK, CritSample, analytical extrema) and
        three interchangeable link representations (AABB, KDOP26, SupportHull).
        This modular interface allows the conservative vs.\ performance trade-off
        to be controlled at runtime without changing the planner or the link
        representation layer.

  \item \textbf{Kinematic-keyed envelope cache (LECT) for multi-query reuse.}
        CCD libraries process each query independently; they do not memoize the
        endpoint or link-envelope evidence across separate planning calls.  LECT
        organises the endpoint evidence as a persistent binary KD-tree over the
        joint-limit box, keyed by robot kinematics rather than by scene geometry,
        so that repeated builds in related scenes can replay stored records for
        unchanged joint intervals.  Schwarzer et al.\ address multi-query reuse
        at the path-segment level, but their cached objects are collision outcomes
        for specific path segments, not kinematic envelope records that survive
        scene changes.

  \item \textbf{Integration into a box-forest planner for repeated queries.}
        The endpoint-to-link envelope map serves an RRT-inspired adaptive
        cell-decomposition planner that grows, consolidates, and reuses a forest
        of validated C-space boxes.  This planning layer---with its sampling
        schedule, frontier-expansion logic, consolidation pass, and corridor-query
        protocol---has no counterpart in the CCD literature above.
\end{enumerate}

\paragraph{Corrected novelty statement.}
The abstract and contribution list in the main paper should be revised to read
(proposed replacement):
\begin{quote}
  \textit{``The core geometric contribution is a modular endpoint-to-link interval
  envelope map for C-space box validation.  The two underlying geometric facts---
  capsule Minkowski decomposition and convex-hull generator bounding of swept
  volumes---are classical and are exploited in CCD libraries (PQP, GJK, FCL) and
  articulated-body CCD work (Redon et al.\ 2004/2005, Zhang et al.\ 2007).
  The novelty in RBF is the combined application to C-space box-growth planning:
  a switchable multi-source endpoint-envelope interface, compact per-representation
  storage (AABB, KDOP26, SupportHull), and kinematic-keyed LECT reuse across
  multi-query builds, none of which is addressed in the prior swept-volume or CCD
  literature.''}
\end{quote}

% ===========================================================================
\section{Replacement Text for Related Work (Section II.B of the Main Paper)}
\label{sec:replacement_text}
% ===========================================================================

The following text replaces the single subsection
\textbf{``II.B Envelope Geometry, Interval Methods, and Collision Checking''}
in the main paper with three subsections, and adds corresponding BibTeX entries.
The revised text explicitly acknowledges the prior pipeline in Redon et al.\
and scopes the RBF claim accordingly.

\bigskip
\hrule
\bigskip

\subsubsection*{II.B-1\quad Swept Volume Approximation and Capsule Geometry}

Computing the swept volume of a moving rigid body has been studied extensively
in computational geometry and CAM communities~\cite{blackmore1992swept,abdel2004swept}.
The exact swept volume of a convex polytope $P$ over a rigid motion is generally
non-convex and expensive to compute; practical methods therefore construct
conservative outer approximations.  For a convex body
$P = \operatorname{conv}(v_1,\dots,v_m)$, bounding each generator trajectory
by an outer set $\widehat{\Gamma}_i$ and forming their convex hull yields
\[
  \mathcal{S}_I(P) \subseteq \operatorname{conv}\!\Bigl(\bigcup_i \widehat{\Gamma}_i(I)\Bigr).
\]
OBB-tree construction uses the same generator-bounding idea for mesh
primitives~\cite{gottschalk1996obbtree}.

The Minkowski sum structure of capsule geometry is equally standard.  A capsule
$C = S \oplus B_2(r)$ is the primitive of PQP, whose swept-sphere volume
construction explicitly decomposes a convex body as a hull plus a fixed-radius
offset~\cite{larsen1999fast}.  The GJK algorithm exploits the same
decomposition~\cite{gilbert1988fast,bergen1999gjk}, and FCL generalises it to
a library of swept-sphere primitives~\cite{pan2012fcl}.

\subsubsection*{II.B-2\quad Interval Forward Kinematics and Endpoint Bounding
  for Articulated Bodies}

Interval arithmetic~\cite{moore1966interval,moore2009introduction,jaulin2001applied}
provides inclusion-monotone enclosures of real-valued functions over parameter
intervals.  Snyder applied interval arithmetic to parametric geometry to compute
conservative bounding boxes for curves and surfaces without mesh
sampling~\cite{snyder1992interval}.  Merlet adapted DH-chain interval
propagation to certified workspace analysis for parallel and serial
manipulators~\cite{merlet2004solving,merlet2009interval}.

These ideas meet in the \emph{continuous collision detection} (CCD) literature for
articulated models.  Redon, Kim, Lin, and Manocha~\cite{redon2004continuous_avatar,
redon2005fast} model each robot link as a line-swept sphere (LSS, i.e.~capsule)
and propagate interval DH transforms along the kinematic chain to obtain
\emph{interval vectors bounding the two LSS endpoint positions} over the motion
time step; the swept capsule is then bounded by padding the endpoint bounding
box with the link radius.  This is explicitly the two-step pipeline of
(interval FK endpoint bounding) $\to$ (Minkowski radius lift) that is the geometric
core of \rbf{}.  Zhang \etal{}~\cite{zhang2007taylor} replace interval arithmetic
with Taylor models in the same pipeline to reduce wrapping errors.  Schwarzer
\etal{}~\cite{schwarzer2005adaptive} address a multi-query motivation similar to
\rbf{}'s: they cache collision evidence for path segments across multiple planning
calls, though at the path-segment level rather than as C-space volumetric records.

\rbf{} uses these same geometric facts and a structurally similar endpoint
bounding step.  The differences from prior CCD work that motivate the paper's
contribution claim are: (a)~the parameter domain is a C-space joint-box interval
rather than a trajectory time step; (b)~the endpoint source is a switchable
interface (IFK, CritSample, analytical extrema) rather than a fixed algorithm;
(c)~compact per-representation link records (AABB, KDOP26, SupportHull) are
selected for staged filtering during box growth; and (d)~the endpoint and link
envelope records are organised in LECT, a persistent kinematic-keyed cache that
reuses evidence across multi-query builds in related scenes---a reuse level not
addressed by CCD libraries or path-segment caches.

\subsubsection*{II.B-3\quad Multi-Query Reuse and the C-Space Box Planning Context}

This paper's use of link interval envelopes differs from CCD in its \emph{planning
purpose}.  CCD asks whether a specific given trajectory collides; it evaluates
a single path and produces a yes/no answer or a time-of-impact.  \rbf{} instead
asks whether \emph{any} configuration in a C-space joint-box is collision-free;
it classifies a volumetric region and retains the validated region for later
reuse.  This shifts the design objectives from tight single-query bounding to
modular, cached, multi-query region classification.  The box-forest planner,
LECT tree, corridor query, and final-audit reporting protocol are consequences
of this C-space-region view of the endpoint-to-link envelope map, and they have
no direct counterpart in the CCD literature above.

\bigskip
\hrule
\bigskip

% ===========================================================================
\section{New BibTeX Entries Required}
\label{sec:new_bibtex}
% ===========================================================================

The following entries should be added to \texttt{references.bib}:

\begin{verbatim}
@techreport{larsen1999fast,
  author      = {Larsen, Eric and Gottschalk, Stefan and Lin, Ming C.
                 and Manocha, Dinesh},
  title       = {Fast Proximity Queries with Swept Sphere Volumes},
  institution = {University of North Carolina at Chapel Hill},
  number      = {TR99-018},
  year        = {1999}
}

@article{bergen1999gjk,
  author  = {Bergen, Gino van den},
  title   = {A Fast and Robust {GJK} Implementation for Collision Detection
             of Convex Objects},
  journal = {J. Graph. Tools},
  volume  = {4},
  number  = {2},
  pages   = {7--25},
  year    = {1999}
}

@inproceedings{snyder1992interval,
  author    = {Snyder, John M.},
  title     = {Interval Analysis for Computer Graphics},
  booktitle = {Proceedings of {ACM SIGGRAPH}},
  pages     = {121--130},
  year      = {1992}
}

@article{blackmore1992swept,
  author  = {Blackmore, Denis and Leu, Ming C.},
  title   = {Analysis of swept volume via {Lie} groups and differential
             equations},
  journal = {Int. J. Robot. Res.},
  volume  = {11},
  number  = {6},
  pages   = {516--537},
  year    = {1992}
}

@article{abdel2004swept,
  author  = {Abdel-Malek, Karim and Yang, Jingzhou and Blackmore, Denis
             and Joy, Kenneth},
  title   = {Swept volumes: foundation, perspectives, and applications},
  journal = {Int. J. Shape Model.},
  volume  = {12},
  number  = {1},
  pages   = {87--127},
  year    = {2006}
}

@inproceedings{redon2004continuous_avatar,
  author    = {Redon, St{\'e}phane and Kim, Young J. and Lin, Ming C.
               and Manocha, Dinesh},
  title     = {Interactive and Continuous Collision Detection for Avatars
               in Virtual Environments},
  booktitle = {{IEEE} Virtual Reality Conference},
  pages     = {117--124},
  year      = {2004}
}

@article{redon2005fast,
  author  = {Redon, St{\'e}phane and Lin, Ming C. and Manocha, Dinesh
             and Kim, Young J.},
  title   = {Fast Continuous Collision Detection for Articulated Models},
  journal = {J. Comput. Inf. Sci. Eng.},
  volume  = {5},
  number  = {2},
  pages   = {126--137},
  year    = {2005}
}

@article{zhang2007taylor,
  author  = {Zhang, Xinyu and Redon, St{\'e}phane and Lee, Minkyoung
             and Kim, Young J.},
  title   = {Continuous Collision Detection for Articulated Models Using
             {Taylor} Models and Temporal Culling},
  journal = {ACM Trans. Graph.},
  volume  = {26},
  number  = {3},
  pages   = {15:1--15:10},
  year    = {2007}
}

@article{schwarzer2005adaptive,
  author  = {Schwarzer, Frank and Saha, Mitul and Latombe, Jean-Claude},
  title   = {Adaptive Dynamic Collision Checking for Single and Multiple
             Articulated Robots in Complex Environments},
  journal = {{IEEE} Trans. Robot.},
  volume  = {21},
  number  = {3},
  pages   = {338--353},
  year    = {2005}
}

@article{corrales2011dynamic,
  author  = {Corrales, Juan Antonio and Candelas, Francisco A.
             and Torres, Fernando},
  title   = {Safe human--robot interaction based on dynamic sphere-swept
             line bounding volumes},
  journal = {Robot. Comput.-Integr. Manuf.},
  volume  = {27},
  number  = {1},
  pages   = {177--185},
  year    = {2011}
}

@book{moore1966interval,
  author    = {Moore, Ramon E.},
  title     = {Interval Analysis},
  publisher = {Prentice-Hall},
  address   = {Englewood Cliffs, NJ},
  year      = {1966}
}
\end{verbatim}

% ===========================================================================
\section{Summary Answer: Has Prior Work Used Endpoint Interval Envelopes\\
  to Compute Link Interval Envelopes for Manipulation?}
\label{sec:summary_answer}
% ===========================================================================

\textbf{Yes, for continuous collision detection (CCD) of articulated bodies.}
The specific two-step pipeline:
\begin{center}
  \emph{interval FK endpoint bounding}
  $\;\longrightarrow\;$
  \emph{endpoint-interval AABB $\oplus$ link radius = link swept-capsule bound}
\end{center}
is \textbf{explicitly present} in:
\begin{itemize}[nosep]
  \item Redon et al., ``Interactive and CCD for Avatars,'' IEEE VR 2004
        \cite{redon2004continuous_avatar} --- avatars/robots as LSS, interval
        endpoint vectors, swept-capsule bounds.
  \item Redon, Lin, Manocha, Kim, ``Fast CCD for Articulated Models,'' JCISE 2005
        \cite{redon2005fast} --- general robot arms, same pipeline, $>240$ citations.
  \item Zhang, Redon, Lee, Kim, ``CCD using Taylor Models,'' ACM TOG 2007
        \cite{zhang2007taylor} --- Taylor-model replacement of IFA, same endpoint-to-link structure.
\end{itemize}

\textbf{No prior work has used this pipeline for C-space box classification
or multi-query planning.}  The complete absence of results on applying
endpoint interval envelopes to C-space joint-box validation and box-forest
planning confirms that the planning-systems contribution of RBF is genuine,
but it requires that the paper \emph{explicitly cite and discuss the CCD
precursors above} and \emph{reframe its novelty claim} from ``first use of
these geometric facts'' to ``first use of these geometric facts for
C-space box growth and multi-query reuse.''

The reviewer concern is therefore \emph{valid}: the paper must acknowledge
Redon et al.\ and Zhang et al.\ as direct predecessors of the endpoint
bounding step, and the contribution must be restated accordingly.
The replacement text in \cref{sec:replacement_text} provides the corrected
Related Work subsections; the corrected contribution item is given in
\cref{sec:novelty_boundary}.

% ===========================================================================
% Bibliography (use the same references.bib from the main paper,
% plus the new entries above)
% ===========================================================================
\bibliographystyle{IEEEtran}

% Inline references for standalone compilation (subset needed for this document)
\begin{thebibliography}{99}

\bibitem{larsen1999fast}
E.~Larsen, S.~Gottschalk, M.~C.~Lin, and D.~Manocha,
``Fast Proximity Queries with Swept Sphere Volumes,''
Technical Report TR99-018, University of North Carolina at Chapel Hill, 1999.

\bibitem{gilbert1988fast}
E.~G.~Gilbert, D.~W.~Johnson, and S.~S.~Keerthi,
``A fast procedure for computing the distance between complex objects in
three-dimensional space,''
\textit{IEEE J. Robot. Autom.}, vol.~4, no.~2, pp.~193--203, 1988.

\bibitem{bergen1999gjk}
G.~van den Bergen,
``A Fast and Robust GJK Implementation for Collision Detection of Convex Objects,''
\textit{J. Graph. Tools}, vol.~4, no.~2, pp.~7--25, 1999.

\bibitem{pan2012fcl}
J.~Pan, S.~Chitta, and D.~Manocha,
``FCL: A general purpose library for collision and proximity queries,''
in \textit{Proc. IEEE ICRA}, 2012, pp.~3859--3866.

\bibitem{gottschalk1996obbtree}
S.~Gottschalk, M.~C.~Lin, and D.~Manocha,
``OBBTree: A hierarchical structure for rapid interference detection,''
in \textit{Proc. ACM SIGGRAPH}, 1996.

\bibitem{blackmore1992swept}
D.~Blackmore and M.~C.~Leu,
``Analysis of swept volume via Lie groups and differential equations,''
\textit{Int. J. Robot. Res.}, vol.~11, no.~6, pp.~516--537, 1992.

\bibitem{abdel2004swept}
K.~Abdel-Malek, J.~Yang, D.~Blackmore, and K.~Joy,
``Swept volumes: foundation, perspectives, and applications,''
\textit{Int. J. Shape Model.}, vol.~12, no.~1, pp.~87--127, 2006.

\bibitem{redon2004continuous_avatar}
S.~Redon, Y.~J.~Kim, M.~C.~Lin, and D.~Manocha,
``Interactive and continuous collision detection for avatars in virtual environments,''
in \textit{Proc. IEEE Virtual Reality Conf.}, 2004, pp.~117--124.

\bibitem{redon2005fast}
S.~Redon, M.~C.~Lin, D.~Manocha, and Y.~J.~Kim,
``Fast continuous collision detection for articulated models,''
\textit{J. Comput. Inf. Sci. Eng.}, vol.~5, no.~2, pp.~126--137, 2005.

\bibitem{zhang2007taylor}
X.~Zhang, S.~Redon, M.~Lee, and Y.~J.~Kim,
``Continuous collision detection for articulated models using Taylor models
and temporal culling,''
\textit{ACM Trans. Graph.}, vol.~26, no.~3, pp.~15:1--15:10, 2007.

\bibitem{schwarzer2005adaptive}
F.~Schwarzer, M.~Saha, and J.-C.~Latombe,
``Adaptive dynamic collision checking for single and multiple articulated robots
in complex environments,''
\textit{IEEE Trans. Robot.}, vol.~21, no.~3, pp.~338--353, 2005.

\bibitem{corrales2011dynamic}
J.~A.~Corrales, F.~A.~Candelas, and F.~Torres,
``Safe human--robot interaction based on dynamic sphere-swept line bounding volumes,''
\textit{Robot. Comput.-Integr. Manuf.}, vol.~27, no.~1, pp.~177--185, 2011.

\bibitem{moore1966interval}
R.~E.~Moore,
\textit{Interval Analysis}.
Englewood Cliffs, NJ: Prentice-Hall, 1966.

\bibitem{moore2009introduction}
R.~E.~Moore, R.~B.~Kearfott, and M.~J.~Cloud,
\textit{Introduction to Interval Analysis}.
Philadelphia, PA: SIAM, 2009.

\bibitem{jaulin2001applied}
L.~Jaulin, M.~Kieffer, O.~Didrit, and E.~Walter,
\textit{Applied Interval Analysis}.
London: Springer, 2001.

\bibitem{snyder1992interval}
J.~M.~Snyder,
``Interval analysis for computer graphics,''
in \textit{Proc. ACM SIGGRAPH}, 1992, pp.~121--130.

\bibitem{merlet2004solving}
J.-P.~Merlet,
``Solving the forward kinematics of a Gough-type parallel manipulator with
interval analysis,''
\textit{Int. J. Robot. Res.}, vol.~23, no.~3, pp.~221--235, 2004.

\bibitem{merlet2009interval}
J.-P.~Merlet,
``Interval analysis for certified numerical solution of problems in robotics,''
\textit{Int. J. Appl. Math. Comput. Sci.}, vol.~19, no.~3, pp.~399--412, 2009.

\end{thebibliography}

\end{document}
